\section*{Anexo 4.1-N --- Correspondencia de componentes lógicos}
\addcontentsline{toc}{section}{Anexo 4.1-N --- Correspondencia de componentes lógicos}
La matriz identifica los doce módulos y los servicios transversales para cotejarlos con el esquema y explicación de solución exigidos en 3.3--3.4 y con su realización física en 4.2. El texto de trabajo disponible del capítulo 3 conserva otra numeración: 3.2.1 contiene alcance funcional de Etapa 1, 3.2.2 infraestructura y 3.3.1 Etapa 2. Estas localizaciones permiten verificar la capacidad existente, pero no sustituyen la adecuación formal de ese capítulo. La correspondencia física es una asignación lógica requerida; su aceptación exige el identificador y la evidencia del componente implantado en el 4.2 integrado.

\begin{tablalafrox}{Cobertura de módulos y realización requerida}{tab:correspondencia-modulos}{F{0.10} F{0.27} F{0.19} Y}{\cab{Módulo} & \cab{Capacidad en capítulo 3} & \cab{Contrato} & \cab{Realización requerida en 4.2}}
M1 & 3.2.1 Recepción & INT-03/06 & Perfil WMS local, base por sitio y ACL ERP. \\
M2 & 3.2.1 Inventario & INT-03/04/12 & WMS por sitio, consolidación central y réplica de lectura de Talca. \\
M3 & 3.2.1 Preventa & INT-01/06 & App Kotlin y API Laravel central. \\
M4 & 3.2.1 y 3.3.1 Ruteo & INT-10 & Servicio de planificación y optimizador separado por contrato. \\
M5 & 3.2.1 Picking y carga & INT-03/07 & HHT, puerta local y WMS con control de guía. \\
M6 & 3.2.1 Entrega y POD & INT-02 & App de reparto, API y almacenamiento de evidencia. \\
M7 & 3.2.1 Rendición; 3.3.1 Cartera & INT-02/06/09 & API, trabajadores Laravel y ACL de cobranza. \\
M8 & 3.2.1 y 3.3.1 Devoluciones & INT-02/03/06 & Registro móvil/local y servicios de conciliación. \\
M9 & 3.2.1 Trazabilidad y frío & INT-03/05 & Bloqueo local, ingesta IoT y consulta de lotes. \\
M10 & 3.3.1 Analítica & Eventos canónicos & Ingesta a lago y servicios de consulta analítica. \\
M11 & 3.3.1 Canal moderno & INT-08 & Reglas Laravel, OpenAS2 y conector por cadena. \\
M12 & 3.3.1 Telemetría & INT-15 & Adaptador de solo lectura y consulta de ruta real. \\
\end{tablalafrox}
{\footnotesize Fuente: capítulo 3 de trabajo, capacidades citadas; catálogo lógico M1--M12 e INT-01--15.\par}

La tabla cubre 12 de 12 módulos, no el 100\% de los componentes del capítulo integrado. Para las capacidades transversales de 3.2.2 se exige correspondencia adicional: presentación y borde con aplicaciones, CloudFront y acceso privado; puerta de enlace con API central y puerta local; eventos con RabbitMQ, shipper y SQS; datos con bases, réplica de lectura y objetos; identidad con Keycloak y verificador local; secretos y auditoría con sus servicios de custodia; observabilidad con OpenTelemetry/ADOT, alarma local y CloudWatch. Notificaciones se traza a INT-11; el motor de optimización de M4 conserva contrato separado y no se presume implementado en PHP por la elección del backend.

El cierre AL-MAP-01 requiere inventariar cada caja de los diagramas 3.3 y 4.1, asignarle identificador estable, contrato, responsable y componente exacto de 4.2, y efectuar la comprobación inversa para detectar infraestructura sin función. La cobertura se calcula como componentes con ambos enlaces verificados dividido por el inventario completo; una referencia genérica a 4.2 no cuenta como enlace verificado. No se declara 100\% hasta integrar las versiones aprobadas de los tres apartados.

